% =====================================================================
% overleaf/09_conclusion.tex — LaTeX rendering of en/09_conclusion.md draft v0 (17 September 2026)
% Dernière mise à jour : 2026-09-17
% Chapter 9 of the paper, a \section of its own, to be inserted right after overleaf/08_limits_and_hybrid.tex.
% This file DEFINES \label{sec:conclusion}.
% TWO THINGS TO TAKE UP BEFORE SUBMISSION, inherited from the French draft and NOT repaired by the rendering:
%   (1) The "Tier 1 / 2 / 3" vocabulary, five occurrences. The written chapters replace it with the SILICA
%       certification steps, exploratory / robust / transferable. Align before submission.
%   (2) The cross-references, which followed the old manuscript numbering. They are rewritten here as \ref to
%       the labels the other .tex files define, so this point is closed ON THE LATEX SIDE only; the Markdown
%       master still carries hard section numbers.
% Citations: bintareaf2026silica and the Baronchelli entries are in references.bib. Check the exact keys against
%   article/CITATIONS.md before compiling: the Markdown writes "Baronchelli, 2026" where the bibliography has
%   baronchelli2025emergent and baronchelli2026groupsize, and the right one for THIS claim is not settled.
% Preamble additions:
%   \newcommand{\todo}[1]{\textbf{[#1]}}
% Structure: five numbered lessons, as an enumerate. No figure, no table.
% Comments "% source:" give the repository file a plain-text figure is recomputed from.
% Derived file: regenerate after every validated pass of en/09_conclusion.md.
% =====================================================================


\section{Conclusion}
\label{sec:conclusion}

\begin{enumerate}

\item \textbf{No LLM for static mass prediction (the Tier 3 ceiling).} In direct agreement with the results of the SILICA benchmark \citep{bintareaf2026silica}, LLM agents fail to reproduce faithfully the empirical distribution of human mobility without heavy statistical calibration. Supervised trees are free in tokens and four times more faithful in argmax: 7.30 points of L1 error against 29.81. % source: scripts/progedo_logit/mode_choice_policy_metrics.json — test.mode_shares.l1_argmax = 0.0730 over 13,045 trips; scripts/synthesis/avancement.yaml — run 2026-08-24_17_34, global L1 29.81 pt

Two reservations attach to that figure. The two errors are not measured on the same substrate: 7.30 comes from the sealed test set of the survey, 29.81 from a simulation run, and the factor of four assumes that they compare term for term, which remains to be established. The speed ratio announced by an earlier version of this work, about 2{,}700 times, has no source in the repository and is withdrawn pending measurement.

\item \textbf{The value of the LLM lies in context dependence and adaptation (Tier 2 validity).} In line with the thesis of Baronchelli, the interest of LLMs does not lie in their status as a statistical human proxy, but in their emergent adaptive behaviour: untabulated dependence, local press, disgust, crowd density, pleasantness; and non-independent dependence, hysteresis at $J+1$ with short-term memory $\mathcal{M}_t$, intra-household arbitration, spatial chain of vehicles. Everywhere else, a tabular model does better and costs less.

\item \textbf{A lesson about measurement, transferable beyond this case.} A variant at 93.4\,\% accuracy on the survey produced the \emph{worst} score in simulation, 9.28 against 7.40 of composite on the ``all decisions'' reading, the distance reconstructed from the declared duration containing the retained mode. Any evaluation of a generative agent must be scored where the model serves, not where it is easy to score.

The ordering was checked under the two readings of the composite, the three published figures being first reproduced identically. Excluding single itineraries, production moves to 10.07 and the variant to 11.31: the variant remains the worst, but the margin falls from 1.88 to 1.24 points. Citing these two numbers without naming the reading is inaccurate. % source: docs/traces/2026-09-12_11-30_ticket047_ordre_des_composites_conclusion/ — run experiments/archive/2026-08-27_17_57, 3,249 decisions of which 665 with a single itinerary

\item \textbf{The hybrid cascade architecture answers the fundamental dichotomy.} Reconcile Baronchelli's two questions by a division of labour: give 90\,\% of the nominal flow to supervised statistical calibration (Tier 3), and reserve generative LLM reasoning for the 10\,\% of complex situations, disruptions and contextual ruptures (Tier 2). Section~\ref{sec:cascade} gives the three stages and says which share is measured and which is not.

\item \textbf{Cost and reproducibility decide where fidelity does not.} A simulated day requires 5{,}040 seconds from the \texttt{gemini-3.5-flash-lite} decider under the expert prompt, against 35 seconds from the kernel regression and from gradient boosting for the same 3{,}154 decisions. The tabular methods return exactly the same probabilities from one run to the next, where the agent changes one in seven. As long as the question asked is the modal split of a territory on an ordinary day, the tabular method therefore remains the appropriate tool: two orders of magnitude cheaper, perfectly reproducible, and not separable from the best agent on the two readings of Section~\ref{sec:agreement}.

Reach works the other way, and it is the only one of the three quantities that does: the tabular methods answer only for the variables they were taught, where the agent receives an instruction in natural language. What Section~\ref{sec:results} establishes is not that a generative agent does better, but that a qualitative instruction free of any numerical threshold brings it within reach of a model that has read 39{,}203 survey trips, and that it installs in it a distance elasticity the model does not have. The interest of the agent begins where the survey stops, and Section~\ref{sec:untabulated} exploits that on two regimes no survey tabulates.

% source: synthese.json, compteurs.duree_s — gemini-3.5 under the minimal prompt (corrected set) 5,613 s for 2,506 solicitations,
%   under prompt_expert_05 (earlier set) 5,040 s for 2,108 solicitations; lgbm 34.8 s and klr 35.2 s on the corrected set,
%   random forest 266 s. The prompt_expert_05 arm of the corrected set was resumed after a quota ran out, 3,019 of its 3,154
%   decisions being re-served: its duration of 540 s measures a resumption, not a full day, and is not a cost reference.
%   The duration of a language-model decider is dominated by the rate limiting of the gateway, not by computation: the ratio
%   measures an operating cost, not a complexity. Reproducibility: indirect measurement, 86.1 % agreement on the most likely
%   mode between two runs of the same prompt on two successive cohorts; identical replay is ticket 073 axis 1.

\end{enumerate}
